% Options for packages loaded elsewhere
\PassOptionsToPackage{unicode}{hyperref}
\PassOptionsToPackage{hyphens}{url}
\PassOptionsToPackage{dvipsnames,svgnames*,x11names*}{xcolor}
%
\documentclass[
  11pt,
]{article}
\usepackage{amsmath,amssymb}
\usepackage{lmodern}
\usepackage{iftex}
\ifPDFTeX
  \usepackage[T1]{fontenc}
  \usepackage[utf8]{inputenc}
  \usepackage{textcomp} % provide euro and other symbols
\else % if luatex or xetex
  \usepackage{unicode-math}
  \defaultfontfeatures{Scale=MatchLowercase}
  \defaultfontfeatures[\rmfamily]{Ligatures=TeX,Scale=1}
  \setmainfont[]{Helvetica}
\fi
% Use upquote if available, for straight quotes in verbatim environments
\IfFileExists{upquote.sty}{\usepackage{upquote}}{}
\IfFileExists{microtype.sty}{% use microtype if available
  \usepackage[]{microtype}
  \UseMicrotypeSet[protrusion]{basicmath} % disable protrusion for tt fonts
}{}
\makeatletter
\@ifundefined{KOMAClassName}{% if non-KOMA class
  \IfFileExists{parskip.sty}{%
    \usepackage{parskip}
  }{% else
    \setlength{\parindent}{0pt}
    \setlength{\parskip}{6pt plus 2pt minus 1pt}}
}{% if KOMA class
  \KOMAoptions{parskip=half}}
\makeatother
\usepackage{xcolor}
\IfFileExists{xurl.sty}{\usepackage{xurl}}{} % add URL line breaks if available
\IfFileExists{bookmark.sty}{\usepackage{bookmark}}{\usepackage{hyperref}}
\hypersetup{
  pdftitle={Real-unit-calibrated simulations close the sim-to-real gap for CNN Einstein-radius regression, with self-diagnosed confidence},
  pdfauthor={Nurkyz Zydyrysova},
  colorlinks=true,
  linkcolor={Maroon},
  filecolor={Maroon},
  citecolor={Blue},
  urlcolor={Blue},
  pdfcreator={LaTeX via pandoc}}
\urlstyle{same} % disable monospaced font for URLs
\usepackage[margin=1in]{geometry}
\usepackage{longtable,booktabs,array}
\usepackage{calc} % for calculating minipage widths
% Correct order of tables after \paragraph or \subparagraph
\usepackage{etoolbox}
\makeatletter
\patchcmd\longtable{\par}{\if@noskipsec\mbox{}\fi\par}{}{}
\makeatother
% Allow footnotes in longtable head/foot
\IfFileExists{footnotehyper.sty}{\usepackage{footnotehyper}}{\usepackage{footnote}}
\makesavenoteenv{longtable}
\setlength{\emergencystretch}{3em} % prevent overfull lines
\providecommand{\tightlist}{%
  \setlength{\itemsep}{0pt}\setlength{\parskip}{0pt}}
\setcounter{secnumdepth}{-\maxdimen} % remove section numbering
\ifLuaTeX
  \usepackage{selnolig}  % disable illegal ligatures
\fi

\title{Real-unit-calibrated simulations close the sim-to-real gap for
CNN Einstein-radius regression, with self-diagnosed confidence}
\author{Nurkyz Zydyrysova}
\date{Draft v0.1}

\begin{document}
\maketitle

\hypertarget{abstract}{%
\subsection{Abstract}\label{abstract}}

\begin{itemize}
\tightlist
\item
  Task: CNN predicts Einstein radius \(\theta_E\) from a single HST
  ACS/WFC F814W lensed image.
\item
  Prior CNN trained on simplified simulations: good in-distribution,
  fails on real lenses (\(R^2 < 0\), prior-pull).
\item
  Fix: paltas/lenstronomy + real COSMOS sources + real STScI
  focus-diverse ePSF + real empty-cutout backdrops + noise calibrated to
  the real SLACS sky-RMS distribution (``hybrid'' training images).
\item
  Result: \(R^2\) crosses negative to positive on 62 SLACS + 40 S4TM
  (frozen, never trained on).
\item
  New addition: \((\mu, \log\sigma^2)\) uncertainty head --- \(\sigma\)
  correlates with real error (\(\rho = +0.71\), SLACS); confidence-gated
  subset meets the target success bar outright.
\item
  Positioning vs.~Cao et al.~2025 {[}cite{]} (same lenses, same
  \(b_{\rm SIE}\) ground truth, conventional pixel-based modeling,
  \textasciitilde3 min/lens): CNN matches on the confident subset at
  roughly \(10^5\) times the inference speed.
\end{itemize}

\hypertarget{introduction}{%
\subsection{1. Introduction}\label{introduction}}

\begin{itemize}
\tightlist
\item
  Strong-lensing \(\theta_E\) is a standard mass-scale probe; upcoming
  surveys (LSST, Euclid, Roman) are expected to discover order \(10^5\)
  lenses --- amortized CNN inference is the only tractable option at
  that scale.
\item
  Prior CNN work:

  \begin{itemize}
  \tightlist
  \item
    Hezaveh, Perreault Levasseur \& Marshall 2017, \emph{Nature}
    {[}cite{]} --- first CNN regression of SIE parameters from HST-like
    simulations.
  \item
    Perreault Levasseur, Hezaveh \& Wechsler 2017, \emph{ApJL}
    {[}cite{]} --- CNN uncertainty estimation for lens parameters.
  \item
    Gawade et al.~2025, arXiv:2404.18897 --- ground-based HSC CNN,
    \textasciitilde10-20\% fractional error.
  \item
    {[}cite, verify not duplicate of above{]} arXiv:1911.06341 --- gri
    wide-field CNN, \(\theta_E\) to \textasciitilde10-15\%.
  \item
    Cao et al.~2025, arXiv:2503.08586 {[}cite{]} --- not a CNN;
    automated pixel-based modeling + nested sampling on 63 grade-A SLACS
    lenses, \(\lesssim 5\%\) deviation, \textasciitilde10\% catastrophic
    failure rate, \textasciitilde3 min/lens. Note their own discussion
    flags ML priors as a path to reducing failures.
  \end{itemize}
\item
  Open problem this paper addresses: CNNs trained on simplified
  simulations show a sim-to-real generalization gap on real telescope
  data.
\item
  Contributions:

  \begin{enumerate}
  \def\labelenumi{\arabic{enumi}.}
  \tightlist
  \item
    Diagnose the prior (pre-paltas) model's real-lens failure as
    \textbf{prior-pull} --- systemic regression toward the training-set
    mean, not an architectural defect.
  \item
    A physically-calibrated ``hybrid'' simulator: paltas/lenstronomy +
    real COSMOS sources + real focus-diverse ePSF (STScI ISR 2018-08 /
    ISR 2023-06 {[}cite{]}) + real empty-field backdrops + noise matched
    to the real sky-RMS distribution --- no hand-tuned flux
    normalization.
  \item
    \(R^2\) crosses negative to positive on the frozen real benchmark;
    the residual gap is quantified via the prediction-vs-truth
    regression slope.
  \item
    An uncertainty head whose \(\sigma\) is domain-aware and predictive
    of real-world failure; gating on it isolates a subset that meets the
    target success bar.
  \end{enumerate}
\end{itemize}

\hypertarget{data}{%
\subsection{2. Data}\label{data}}

\hypertarget{real-benchmark-frozen-never-used-in-training}{%
\subsubsection{2.1 Real benchmark (frozen, never used in
training)}\label{real-benchmark-frozen-never-used-in-training}}

\begin{itemize}
\tightlist
\item
  \(N = 62\) SLACS + 40 S4TM {[}cite S4TM{]}, HST ACS/WFC F814W.
\item
  Ground truth: Bolton et al.~2008 SIE \(b_{\rm SIE}\) {[}cite{]}.
\item
  One SLACS lens (J0955+0101) excluded --- corrupted cutout.
\item
  Cutouts: 128x128 px at 0.05 arcsec/px (6.4 arcsec field of view), MAST
  drizzled products.
\end{itemize}

\hypertarget{training-simulator-hybrid}{%
\subsubsection{2.2 Training simulator
(``hybrid'')}\label{training-simulator-hybrid}}

\begin{itemize}
\tightlist
\item
  Base: paltas {[}cite{]} on lenstronomy; PEMD+shear deflector;
  COSMOS\_23.5 source catalog (real, GREAT3-vetted galaxies).
\item
  \(\theta_E \sim U[0.45, 2.30]\) arcsec, flat.
\item
  Lens light: joint empirical (magnitude, \(R_{\rm Sersic}\)) prior,
  drawn from 63 (magnitude, \(R_e\)) pairs measured directly from the
  real SLACS cutouts, with jitter.
\item
  PSF: STScI focus-diverse empirical PSFs (real HST optics), 88 distinct
  kernels from real archival exposures, randomly rotated per generation
  shard.
\item
  Backdrop: noiseless simulated render + a real empty COSMOS cutout
  (real correlated noise and field neighbors) + Gaussian top-up noise
  drawn per-image from the real SLACS sky-RMS distribution.
\end{itemize}

\textbf{Table 1. Dataset realism gate (sim vs.~real, identical estimator
both sides).}

\begin{longtable}[]{@{}llll@{}}
\toprule
Metric & Simulated & Real & Target \\
\midrule
\endhead
Sky RMS ratio (sim/real) & 1.14 & --- & 0.8 - 1.25 \\
Lens peak / sky & 723 & 492 - 1798 (16-84\%) & within real band \\
\(\theta_E\) range & {[}0.45, 2.30{]} & {[}0.54, 1.78{]} & benchmark
inside interior \\
\bottomrule
\end{longtable}

\begin{itemize}
\tightlist
\item
  100,000 training / 5,000 validation images; validation uses seed- and
  PSF-kernel-disjoint generation.
\end{itemize}

\textbf{Figure 1.} Real vs.~simulated cutouts under an identical
stretch. \emph{{[}insert figure{]}}

\hypertarget{current-limitations-of-the-simulator}{%
\subsubsection{2.3 Current limitations of the
simulator}\label{current-limitations-of-the-simulator}}

\begin{itemize}
\tightlist
\item
  Lens light is a parametric Sersic profile, not a real elliptical
  galaxy image (real deflectors show disky/inclined morphology, dust,
  tidal features).
\item
  Empty-field backdrops drawn from a single COSMOS tile (1,317 unique
  cutouts, reused across the training set); a larger backdrop pool is in
  preparation.
\item
  Mass model is SIE/PEMD only; no hydrodynamical-simulation mass
  structure yet (see Section 6).
\end{itemize}

\hypertarget{model-and-training}{%
\subsection{3. Model and training}\label{model-and-training}}

\begin{itemize}
\tightlist
\item
  Architecture: scale-conditioned ResNet (four residual stages, global
  average pooling, scale-scalar concatenation, MLP head).
\item
  Input: single-band, asinh-normalized image (per-image mean/std after
  arcsinh compression).
\item
  Loss: point estimate (Huber) in an initial version;
  \((\mu, \log\sigma^2)\) Gaussian negative log-likelihood in the
  reported model (uncertainty head).
\item
  Optimizer: Adam, cosine learning-rate schedule, 40 epochs, batch size
  64; flip/rotation augmentation (\(\theta_E\) is invariant under all
  eight dihedral transforms).
\item
  Model selection: simulation validation split only; the real benchmark
  is evaluated at fixed checkpoints, never used for selection.
\item
  Test-time augmentation: prediction averaged over the eight dihedral
  views; total predicted uncertainty combines the mean aleatoric
  variance with the across-view (epistemic) variance.
\end{itemize}

\hypertarget{results}{%
\subsection{4. Results}\label{results}}

\hypertarget{sim-to-real-failure-diagnosis-baseline-pre-hybrid-pipeline}{%
\subsubsection{4.1 Sim-to-real failure diagnosis (baseline,
pre-hybrid-pipeline)}\label{sim-to-real-failure-diagnosis-baseline-pre-hybrid-pipeline}}

\textbf{Table 2. Baseline model on the frozen real benchmark.}

\begin{longtable}[]{@{}llllll@{}}
\toprule
Sample & \(N\) & Median frac. error & \(R^2\) & MAE & Failure rate
(\(>15\%\)) \\
\midrule
\endhead
SLACS & 62 & \(-7.9\%\) & \(-1.02\) & 0.274 arcsec & 55\% \\
S4TM & 40 & \(-5.3\%\) & \(-0.53\) & --- & 68\% \\
\bottomrule
\end{longtable}

\begin{itemize}
\tightlist
\item
  Diagnosis: signed error sweeps from positive to negative across true
  \(\theta_E\), crossing zero near a training-mean attractor
  (\textasciitilde0.95-1.0 arcsec) --- prior-pull. Weak per-feature
  correlation with the residuals (all \(|\rho| < 0.3\)) indicates a
  systemic effect rather than an architectural or image-specific one.
\end{itemize}

\hypertarget{hybrid-pipeline-iteration-1}{%
\subsubsection{4.2 Hybrid pipeline, iteration
1}\label{hybrid-pipeline-iteration-1}}

\textbf{Table 3. Hybrid model, iteration 1, on the frozen real benchmark
(bootstrap 95\% CI, 10,000 resamples).}

\begin{longtable}[]{@{}llllll@{}}
\toprule
Sample & \(N\) & Median frac. error (95\% CI) & \(R^2\) & MAE & Failure
rate (95\% CI) \\
\midrule
\endhead
SLACS & 62 & \(-1.1\%\) {[}\(-4.0\), \(+2.2\){]} & \(+0.23\) & 0.139
arcsec & 27\% {[}18, 39{]} \\
S4TM & 40 & \(+5.8\%\) {[}\(+2.1\), \(+16.9\){]} & \(+0.33\) & 0.159
arcsec & 38\% {[}22, 52{]} \\
\bottomrule
\end{longtable}

\begin{itemize}
\tightlist
\item
  \(R^2\) crosses from negative to positive on both samples relative to
  the baseline.
\item
  The residual S4TM bias traces to lenses below the training prior's
  lower bound (true \(\theta_E\) down to 0.54 arcsec against a training
  floor of 0.6 arcsec).
\end{itemize}

\hypertarget{hybrid-pipeline-iteration-2-widened-prior-expanded-psf-diversity-uncertainty-head-test-time-augmentation}{%
\subsubsection{4.3 Hybrid pipeline, iteration 2 (widened prior, expanded
PSF diversity, uncertainty head, test-time
augmentation)}\label{hybrid-pipeline-iteration-2-widened-prior-expanded-psf-diversity-uncertainty-head-test-time-augmentation}}

\textbf{Table 4. Hybrid model, iteration 2, on the frozen real
benchmark.}

\begin{longtable}[]{@{}lllllll@{}}
\toprule
Sample & \(N\) & Median frac. error (95\% CI) & \(R^2\) & MAE & Failure
rate (95\% CI) & Slope \\
\midrule
\endhead
SLACS & 62 & \(-2.6\%\) {[}\(-4.5\), \(-0.1\){]} & \(+0.27\) & 0.129
arcsec & 23\% {[}13, 34{]} & 0.72 \\
S4TM & 40 & \(-1.8\%\) {[}\(-7.9\), \(+1.5\){]} & \(+0.47\) & 0.140
arcsec & 38\% {[}22, 52{]} & 0.71 \\
\bottomrule
\end{longtable}

\begin{itemize}
\tightlist
\item
  Both medians now fall within a \(\pm 5\%\) target band; the
  prediction-vs-truth regression slope (1.0 = no compression) improves
  from 0.62 to 0.72 (SLACS) and 0.58 to 0.71 (S4TM), quantifying a
  reduced but still-present domain-gap compression.
\end{itemize}

\textbf{Figure 2.} Predicted vs.~true \(\theta_E\): baseline model,
hybrid iteration 2, and the Cao et al.~2025 accuracy band.
\emph{{[}insert figure{]}}

\hypertarget{confidence-gated-performance}{%
\subsubsection{4.4 Confidence-gated
performance}\label{confidence-gated-performance}}

\begin{itemize}
\tightlist
\item
  Spearman correlation between predicted \(\sigma/\mu\) and true
  absolute fractional error: \(+0.71\) (SLACS), \(+0.39\) (S4TM) --- the
  model's own uncertainty estimate is predictive of real-world error,
  not only of simulation-validation error.
\item
  \(\sigma/\mu\) is itself domain-aware: a median of roughly 9-10\% on
  real lenses versus roughly 1\% on simulation validation, without any
  explicit domain label at training time.
\end{itemize}

\textbf{Table 5. Confidence-gated subsets.}

\begin{longtable}[]{@{}lllll@{}}
\toprule
Sample & Subset & \(N\) & Median frac. error & Failure rate \\
\midrule
\endhead
SLACS & full & 62 & \(-2.6\%\) & 23\% \\
SLACS & \(\sigma/\mu \le\) median & 31 & \(-1.6\%\) & \textbf{6\%} \\
SLACS & \(\sigma/\mu \le\) 75th pct. & 46 & \(-1.5\%\) & 7\% \\
S4TM & full & 40 & \(-1.8\%\) & 38\% \\
S4TM & \(\sigma/\mu \le\) median & 20 & \(-0.9\%\) & 15\% \\
\bottomrule
\end{longtable}

\begin{itemize}
\tightlist
\item
  The confident half of the SLACS sample meets the target success bar
  outright (\(\pm 5\%\) median, \(\lesssim 10\%\) failure) and is
  competitive with the Cao et al.~2025 failure rate at substantially
  higher inference speed.
\end{itemize}

\hypertarget{planned-but-not-yet-completed}{%
\subsubsection{4.5 Planned but not yet
completed}\label{planned-but-not-yet-completed}}

\begin{itemize}
\tightlist
\item
  Ablation: flat vs.~skewed \(\theta_E\) prior.
\item
  Ablation: benchmark-matched vs.~broad (non-benchmark) PSF pool.
\item
  Ablation: real backdrop vs.~Gaussian noise only.
\item
  Ablation: real empirical PSF vs.~Gaussian/Moffat PSF.
\item
  Matched per-lens comparison against Cao et al.~2025's public per-lens
  results.
\item
  S4TM-specific noise recalibration (current noise model calibrated to
  SLACS sky-RMS only).
\item
  Multi-parameter (ellipticity, centroid) re-validation on the hybrid
  pipeline.
\end{itemize}

\hypertarget{discussion}{%
\subsection{5. Discussion}\label{discussion}}

\begin{itemize}
\tightlist
\item
  The crossing of \(R^2\) from negative to positive is the central
  result: the baseline model did not track individual real lenses; the
  hybrid-trained model does.
\item
  The residual gap is compression (slope \(<1\)) rather than bias
  (median \(\approx 0\)): the model hedges toward the prior mean under
  domain shift, a milder form of the baseline's prior-pull.
\item
  The uncertainty head serves three purposes: a calibration/error-bar
  deliverable, a practical deployment filter, and an ingredient
  (mean-variance estimation) that pairs naturally with unsupervised
  domain adaptation (Section 6.2).
\item
  Standing limitation: the lens light remains parametric. If a fully
  realism-calibrated model still showed \(R^2 < 0\) on real lenses, that
  would motivate escalating to real deflector-cutout injection; this has
  not been triggered, as \(R^2\) is positive.
\end{itemize}

\hypertarget{proposed-extensions}{%
\subsection{6. Proposed extensions}\label{proposed-extensions}}

\hypertarget{illustristng-convergence-maps-within-paltas}{%
\subsubsection{6.1 IllustrisTNG convergence maps within
paltas}\label{illustristng-convergence-maps-within-paltas}}

\begin{itemize}
\tightlist
\item
  Replace or augment the SIE/PEMD deflector with deflection fields
  derived from IllustrisTNG convergence maps, introducing real
  (non-elliptical) mass structure: twists, substructure, boxiness.
\item
  Two risks require checking before use: (i) a label-convention offset
  between SIE \(b_{\rm SIE}\) and the mean-convergence-based
  \(\theta_E\) definition, to be quantified on a TNG subsample before
  mixing label types; (ii) the known IllustrisTNG central-image artifact
  from limited hydrodynamical resolution {[}Bolton 2012; Shu et al.~2016
  --- cite{]}, a new sim-to-real gap that this change would introduce
  and that must be checked for.
\end{itemize}

\hypertarget{sim-to-real-domain-adaptation}{%
\subsubsection{6.2 Sim-to-real domain
adaptation}\label{sim-to-real-domain-adaptation}}

\begin{itemize}
\tightlist
\item
  Prior domain-adaptation work for lens \(\theta_E\) regression exists
  but adapts between two simulated domains, not from simulation to real
  data: Ciprijanovic et al.~{[}cite, arXiv:2311.17238{]} (adversarial
  and maximum-mean-discrepancy adaptation) and {[}cite,
  arXiv:2411.03334{]} (unsupervised domain adaptation combined with a
  mean-variance estimator).
\item
  This leaves open a sim-to-real domain adaptation validated against a
  real spectroscopic benchmark, which this project's frozen 62+40 sample
  with \(b_{\rm SIE}\) ground truth is positioned to provide.
\item
  Any unsupervised adaptation must draw its unlabeled target pool from
  real lenses disjoint from the frozen benchmark, to avoid contaminating
  the evaluation set.
\item
  Sequenced after the ablations in Section 4.5, so that the
  pre-adaptation sim-to-real gap is fully characterized first.
\end{itemize}

\hypertarget{conclusion}{%
\subsection{7. Conclusion}\label{conclusion}}

\begin{itemize}
\tightlist
\item
  {[}to be written after ablations are complete{]}
\end{itemize}

\hypertarget{references}{%
\subsection{References}\label{references}}

\begin{itemize}
\tightlist
\item
  Bolton et al.~2008, SLACS \(b_{\rm SIE}\) catalog {[}cite{]}
\item
  Hezaveh, Perreault Levasseur \& Marshall 2017, \emph{Nature} 548, 555
  {[}cite{]}
\item
  Perreault Levasseur, Hezaveh \& Wechsler 2017, \emph{ApJL} {[}cite{]}
\item
  Gawade et al.~2025, arXiv:2404.18897
\item
  {[}cite{]} arXiv:1911.06341
\item
  Cao et al.~2025, arXiv:2503.08586 {[}cite{]}
\item
  Wagner-Carena et al., paltas {[}cite{]}
\item
  Ciprijanovic et al., arXiv:2311.17238
\item
  {[}cite{]} arXiv:2411.03334
\item
  {[}cite{]} arXiv:2410.16347
\item
  Bolton 2012; Shu et al.~2016 {[}cite{]}
\item
  Wan, Chan, Lange \& Hannuksela, LensFusion {[}cite{]}
\end{itemize}

\end{document}
